%%%%%%%%%%%%%%%%%%%%%%%%%%%%%%%%%%%%%%%%%%%%%%%%%%%%%%%%%%%%%%%%%%%%%%%%%%%
\chapter{Theoretical Framework} 
%may be changed into a Technical Background
%%%%%%%%%%%%%%%%%%%%%%%%%%%%%%%%%%%%%%%%%%%%%%%%%%%%%%%%%%%%%%%%%%%%%%%%%%%

\section{Cognitive Load Theory}
Cognitive load theory (CLT) provides a framework for understanding human cognitive strain and classifies cognitive load into 3 main categories: intrinsic, extraneous, and germane \cite{sweller2019architecture}. Each of these are defined below:
\begin{itemize}
    \item Intrinsic cognitive load is the load induced by the task itself. Intrinsic cognitive load is only affected by two factors: the person and the activity itself. A person’s experience is a key factor in the difficulty induced by a task, because an expert in a task, someone with prior experience with the task or a similar task, receives less cognitive load than a novice, a person with no experience, doing the same task. The difficulty of the task itself affects intrinsic cognitive load, with more difficult tasks increasing cognitive load for the same participant.
\end{itemize}
\begin{itemize}
    \item Extraneous cognitive load is the load induced by the format of information and the environment. Different ways of presenting information induce different cognitive loads. A primary example of this is how visual and auditory media (video) induce lower cognitive load than visual presentation of information alone, a case of the modality effect. The environment also contributes to extraneous cognitive load, where additional stimulation such as noise or discomfort increases cognitive load.
\end{itemize}
\begin{itemize}
    \item Germane cognitive load is the cognitive load spent on learning. Low to moderate cognitive load allows for germane load to occur, whereas cognitive overload does not allow allocation of cognitive resources to allow for learning.
\end{itemize}
These three categories underpin a framework for understanding the factors affecting human cognitive strain. By developing an understanding of human cognitive load during activities, managing cognitive load becomes feasible. As technology and software integrate into the daily lives of people, understanding the implications of using these tools becomes increasingly important. Cognitive load theory can provide a theoretical approach for evaluating human experience while using technology. In this study, CTL will help the researchers identify the effect of different contemporary systems on the cognitive load of neurodivergent and neurotypical users, particularly the intrinsic and extraneous cognitive load associated with the use of software systems themselves.


\section{Pupillometry}
Pupil dilation has been used as a metric for cognitive load in recent years \cite{souchet2022, vilotijevic2022, castanheira2021}. Increased cognitive strain correlates with the dilation of pupils and can be used to map periods of heightened cognitive load, and this approach can be a reliable quantitative metric for evaluating cognitive load. There are other factors that affect pupil size, such as pupil light response, and these generally start at 200 ms, while the onset of task-evoked pupil dilation, such as those in this research, tends to be slow-paced; researchers should observe 2–3 seconds after the onset of stimulation for pupil response \cite{mathot2022}.

\section{Neurodiversity Theory}
Neurodiversity is a cognitive condition that uniquely characterizes an individual’s way of thinking \cite{baumer2021, childmind2025}. This theory highlights the natural variation of the human factor in HCI studies. It must be established that neurodiversity is associated with a unique cognitive ability rather than a deficit \cite{aldersonday2023}, because there exists a trend wherein neurodivergent individuals are kept from co-producing tools and systems that are dedicated to them \cite{burtscher2024, le2024, baumer2021, spiel2022, stefanidi2022}. However, the reason behind why it happens is unknown. Therefore, examining the impacts of neurodiversity on the development of tools, applications, and systems can help create a framework that defines effective approaches and practices grounded in this theory. 
 

\section{Integration of CLT and Neurodiversity}
The interconnection of CLT and neurodiversity is an extremely under-examined area \cite{lecunff2024}. Currently, there is a noticeable lack of frameworks that describe the relationship between these domains, especially in the context of HCI, user experience, and systems design. This extremely novel area of study creates an opportunity for invaluable discoveries and syntheses. This study contributes to the field by proposing a framework that integrates the core concepts of CLT and neurodiversity in relation to the above-mentioned domains. By grounding the framework in both empirical and theoretical findings, this study can establish foundational best practices for future research, development, and implementation of accessible digital tools and environments.
